\originalchapter{概率与物理观测}\label{ch:prob}
\originalsection{微观不确定性与概率密度}
同一温度的气体中, 分子速度并不相同; 记录电阻两端的电压, 也会发现它围绕平均值起伏. 概率描述的对象是重复准备、不同粒子或不同观测时刻构成的样本, 并不是把某一次未知的数值改成确定平均值. 经典统计中的不确定性可来自初始条件未知, 量子统计还包含测量本身的概率规律.

连续随机变量$X$的密度$p_X(x)$满足$p_X(x)\ge0$和$\int p_X(x)\dd x=1$, 区间概率是密度在该区间的积分. $p_X$的单位是$X$单位的倒数, 密度本身可以大于1. 分布函数$P_X(x)=\Pr(X\le x)$没有单位. 离散变量用概率质量$p_n=\Pr(X=x_n)$, 也可统一写为$\sum_n p_n\delta(x-x_n)$, 其中Dirac函数表示单位面积的集中概率. 例如脱附前仍停留在原位的原子与已离开的原子组成混合分布
\[
p(x,t)=\e^{-t/\tau}\delta(x)+(1-\e^{-t/\tau})f(x),\qquad \int f(x)\dd x=1.
\]
$\tau$是平均等待时间, 指数项是尚未离开的概率. 原点处是概率质量, 不能把它当成有限高度的柱子.

观测函数$f(X)$的平均为$\avg{f(X)}=\int f(x)p_X(x)\dd x$. 方差$\Var X=\avg{X^2}-\avg X^2$衡量偏离平均值的平方尺度, 标准差为其平方根. 高斯密度$\exp[-(x-\mu)^2/(2\sigma^2)]/\sqrt{2\pi\sigma^2}$中$\mu$和$\sigma$独立控制中心与宽度. 指数等待时间密度$\e^{-t/\tau}/\tau$在$t\ge0$上成立, 其均值为$\tau$, 方差为$\tau^2$.

若独立事件以单位长度$r$的恒定平均率发生, 将长度$L$分为$M$小段, 每段发生一次的概率为$rL/M+o(M^{-1})$, 发生多次的概率为$o(M^{-1})$. 在$M\to\infty$下, 二项概率给出
\[
 p_n=\lim_{M\to\infty}\binom Mn(rL/M)^n(1-rL/M)^{M-n}
 =\e^{-rL}(rL)^n/n!.
\]
这是Poisson分布, $\avg n=\Var n=rL$. 由$p_0=\e^{-rt}$求第一次事件的等待时间, 又得到指数密度. 每次碰壁以概率$q$从小孔逃逸的原子, 在逃逸前碰壁$n$次的概率则为$q(1-q)^n$, $n\ge0$. 连续时间等待与离散次数等待是两种不同的观测.

\originalsection{条件概率与相关}
两变量的联合密度$p(x,y)$给出一个小矩形中的概率$p(x,y)\dd x\dd y$. 忽略$Y$时需对$y$积分, 得边缘密度$p_X(x)=\int p(x,y)\dd y$. 已知$Y=y$且$p_Y(y)>0$时, 条件密度是$p(x\mid y)=p(x,y)/p_Y(y)$. 连续事件$Y=y$自身通常概率为零, 此式的含义是小区间条件概率的极限, 不是除以一个零概率.

统计独立要求$p(x,y)=p_X(x)p_Y(y)$, 比协方差为零强. 圆盘内均匀分布的$(X,Y)$有零协方差, 但已知$Y=y$便把$X$限制在$[-\sqrt{a^2-y^2},\sqrt{a^2-y^2}]$, 因而不独立. 直接积分给出$p_Y(y)=2\sqrt{a^2-y^2}/(\pi a^2)$, 条件密度为该弦长的倒数.

噪声电压在相隔$\tau$的两时刻可近似用等方差联合高斯表示:
\[
p(v_1,v_2)=\frac{\exp[-(v_1^2-2\rho v_1v_2+v_2^2)/(2\sigma^2(1-\rho^2))]}{2\pi\sigma^2\sqrt{1-\rho^2}},\quad |\rho|<1.
\]
$\rho(\tau)=\avg{v(t)v(t+\tau)}/\sigma^2$是相关系数, 两个边缘均为零均值、方差$\sigma^2$的高斯. 配方得到条件均值$\rho v_2$, 条件方差$\sigma^2(1-\rho^2)$. 因而短时间相关强时, 一次测量有助于预测下一次; $\rho=0$时本模型才给出独立. 极限$\rho\to1$成为集中在$v_1=v_2$上的奇异分布, 不能直接代入普通二维密度.

\originalsection{变量变换与物理通量}
仪器测到的量常常不是原来的随机变量. 若$Y=f(X)$, 安全的起点是$\Pr(Y\le y)=\int_{f(x)\le y}p_X(x)\dd x$. 在每一单调分支上求导, 给出
\[
p_Y(y)=\sum_{x_j:f(x_j)=y}\frac{p_X(x_j)}{|f'(x_j)|},
\]
求和涵盖所有原像, 只在分支可微且导数非零处使用. 例如光强$I=aE^2$, $a>0$, 电场$E$为零均值高斯, 必须同时计入$\pm\sqrt{I/a}$:
\[
p_I(I)=\frac{\e^{-I/(2a\sigma^2)}}{\sqrt{2\pi a\sigma^2 I}},\quad I>0.
\]
它在原点发散但可积, 这并不意味着原点有有限概率质量. 频率与波长满足$\lambda=c/\nu$, 因而$p_\lambda(\lambda)=c\,p_\nu(c/\lambda)/\lambda^2$; 一个谱的峰不能只按变量值取倒数, 必须连同Jacobian一起转换. 两个独立$[0,1]$均匀变量的最大值$Z$满足$P_Z(z)=z^2$, 因而$p_Z=2z$, $\avg Z=2/3$, $\Var Z=1/18$.

穿过平面的粒子比平面附近随机挑选的粒子多一项法向速度权重. 平衡气体速度分量各为方差$\sigma^2=\kb T/m$的高斯. 这里$T$是绝对温度, $m$是粒子质量; $\kb$是Boltzmann（玻尔兹曼）常量, 单位为$\mathrm{J/K}$, 将温度转换成热运动的能量尺度$\kb T$. 从表面向外的通量分布中, 速度大小、极角与方位角分别为
\[
p(v)=\frac{v^3}{2\sigma^4}\e^{-v^2/(2\sigma^2)},\quad
p(\theta)=2\sin\theta\cos\theta,\quad p(\phi)=\frac1{2\pi},
\]
其中$v\ge0$, $0\le\theta\le\pi/2$. 因为$\dd\Omega=\sin\theta\dd\theta\dd\phi$, 单位立体角概率为$\cos\theta/\pi$, 这就是余弦通量律. 横穿相距$d$的平面时路程为$d/\cos\theta$, 所以平均路程是$\int_0^{\pi/2}2d\sin\theta\dd\theta=2d$. 对$v_z=v\cos\theta$积分所有横向速度, 或对分布函数求导, 得
\[
p(v_z)=\frac{v_z}{\sigma^2}\e^{-v_z^2/(2\sigma^2)},\quad v_z>0.
\]
普通气体的速度高斯、普通气体的速率Maxwell分布与逃逸通量的速率分布必须分清.

\originalsection{求和、极限与相对涨落}
两个变量之和$S=X+Y$的密度是$\int p(x,s-x)\dd x$, 独立时成为卷积$\int p_X(x)p_Y(s-x)\dd x$. 有限二阶矩下, 总和方差为各方差加两两协方差的两倍, 只有独立或不相关时才可删去交叉项. 例如前述两次高斯电压之和的方差是$2\sigma^2(1+\rho)$.

原子氢maser中每次在壁面停留的时间若独立且为均值$\tau$的指数变量, $n$次总停留时间有密度
\[
p(t\mid n)=\frac{t^{n-1}}{(n-1)!\tau^n}\e^{-t/\tau},\quad t\ge0,\ n\ge1.
\]
从$n=1$开始, 每次卷积的积分是$\int_0^t u^{n-1}\dd u=t^n/n$, 便归纳得到此式. 把整数阶乘推广到正实参数时, 用Gamma函数
\[\Gamma(a)=\int_0^\infty x^{a-1}\e^{-x}\dd x,\qquad a>0.\]
$x$和$a$在这里无量纲. 原点附近的积分要求$a>0$, 无穷远由指数衰减保证收敛; 分部积分给$\Gamma(a+1)=a\Gamma(a)$, 所以$\Gamma(n)=(n-1)!$, 且Gaussian积分给$\Gamma(1/2)=\sqrt\pi$. 上式停留时间密度中的$(n-1)!$可写成$\Gamma(n)$; 对一般形状参数$a>0$, Gamma密度是$t^{a-1}\e^{-t/\tau}/[\Gamma(a)\tau^a]$, $t\ge0$, $\tau>0$. 同尺度的Gamma密度卷积仍为Gamma密度, 形状参数相加. 这比“求和后总为高斯”准确.

求和的极限可用\textbf{特征函数}描述. 对任意实随机变量$Y$, 定义
\[\varphi_Y(t)=\avg{\e^{\ii tY}};
\quad\text{有密度时}\quad\varphi_Y(t)=\int_{-\infty}^{\infty}\e^{\ii ty}p_Y(y)\dd y.\]
离散变量则把积分换成概率加权和. $t$的单位是$Y$单位的倒数, $\ii^2=-1$; 因为复指数的模恒为一, 即使某些矩不存在, 此平均也存在. 独立的$Y_i$满足$\varphi_{\sum_iY_i}(t)=\prod_i\varphi_{Y_i}(t)$: 指数把和化成乘积, 独立性再把平均拆开. 有限二阶矩允许在$t=0$附近展开, 对中心化$Y=X-\mu$得到$\varphi_Y(t)=1-\sigma^2t^2/2+o(t^2)$.

因此独立同分布、有限非零方差的$X_i$满足中心极限定理: $Z_n=(\sum_iX_i-n\mu)/(\sigma\sqrt n)$的分布函数趋于标准高斯. 对每个固定$t$, 上述乘积给
\[\varphi_{Z_n}(t)=\left[\varphi_Y\!\left(\frac{t}{\sigma\sqrt n}\right)\right]^n
=\left[1-\frac{t^2}{2n}+o(n^{-1})\right]^n\longrightarrow\e^{-t^2/2}.\]
最后一步用$(1+c/n)^n\to\e^c$. Gaussian积分表明右端正是标准高斯的特征函数. 此处使用的连续性定理是: 若特征函数逐点趋于某一分布的特征函数, 且极限在原点连续, 则分布函数在极限分布的连续点收敛. 标准高斯在各点连续, 故得到所述结论. 它不保证密度逐点收敛, 离散变量仍在格点上. 非同分布或有相关的系统要另查Lindeberg条件或相关衰减条件. Cauchy变量没有有限方差, 不在这个定理的适用范围内.

若$\mu\ne0$, 总和平均量级$n$, 标准差量级$\sqrt n$, 相对涨落量级$n^{-1/2}$. 这解释了大量粒子组成的宏观系统为何稳定, 同时也指明了例外: 强相关、临界点或少数项支配时, 这个估计可能失效. 正文后部会把这个概率尺度与可测的热容、压缩率相连.
